
\documentclass{article}
\usepackage{colortbl}
\usepackage{makecell}
\usepackage{multirow}
\usepackage{supertabular}

\begin{document}

\newcounter{utterance}

\centering \large Interaction Transcript for game `chronicle', experiment `battles\_hard', episode 6 with Bentel rockers\_\_qwen3.5{-}4b\_\_Bentel\_iter\_3.
\vspace{24pt}

{ \footnotesize  \setcounter{utterance}{1}
\setlength{\tabcolsep}{0pt}
\begin{supertabular}{c@{$\;$}|p{.15\linewidth}@{}p{.15\linewidth}p{.15\linewidth}p{.15\linewidth}p{.15\linewidth}p{.15\linewidth}}
    \# & $\;$A & \multicolumn{4}{c}{Game Master} & $\;\:$B\\
    \hline

    \theutterance \stepcounter{utterance}  
    & & \multicolumn{4}{p{0.6\linewidth}}{
        \cellcolor[rgb]{0.9,0.9,0.9}{
            \makecell[{{p{\linewidth}}}]{
                \texttt{\tiny{[P1$\langle$GM]}}
                \texttt{You are playing a historical deduction game as the NARRATOR.} \\
\\ 
\texttt{You have been given the name of a historical event, discovery, or milestone. Your goal is to help the DETECTIVE figure out which event you are describing — but you must describe it indirectly, without ever naming it.} \\
\\ 
\texttt{Rules:} \\
\texttt{1. Each round, write exactly ONE sentence describing the event. Add NEW information each round — do not repeat what you wrote before.} \\
\texttt{2. You must NEVER use any of the following forbidden keywords (these would make the event too easy to identify): shi le, xiangguo} \\
\texttt{3. You must NEVER state the event name directly or any close variant of it.} \\
\texttt{4. You may freely describe causes, consequences, the broader context, the sequence of events, and the significance — but only using general, non{-}identifying language.} \\
\\ 
\texttt{Format your response EXACTLY as follows (required):} \\
\texttt{PARAGRAPH: <exactly one sentence describing the event without forbidden keywords>} \\
\\ 
\texttt{The event you must describe (without naming it) is: Battle of Xiangguo (312)} \\
            }
        }
    }
    & \\ \\

    \theutterance \stepcounter{utterance}  
    & \multicolumn{4}{p{0.6\linewidth}}{
        \cellcolor[rgb]{0.9,0.9,0.9}{
            \makecell[{{p{\linewidth}}}]{
                \texttt{\tiny{[P1$\rangle$GM]}}
                \texttt{PARAGRAPH: In the year 312, a fierce confrontation erupted between the forces of Cao Shuang and Sima Yi at the capital, marking a decisive turning point in the struggle for control of the Wei state.} \\
            }
        }
    }
    & & \\ \\

    \theutterance \stepcounter{utterance}  
    & & \multicolumn{4}{p{0.6\linewidth}}{
        \cellcolor[rgb]{0.9,0.9,0.9}{
            \makecell[{{p{\linewidth}}}]{
                \texttt{\tiny{[GM$\rangle$P2]}}
                \texttt{You are playing a historical deduction game as the DETECTIVE.} \\
\\ 
\texttt{A Narrator will describe a historical event, discovery, or milestone to you over several rounds. Each round you receive a new sentence of description that builds on what came before. Your job is to figure out which specific event is being described.} \\
\\ 
\texttt{After EACH new sentence, respond briefly. Keep your entire response under 100 words. Use this exact format (both sections are required):} \\
\\ 
\texttt{ANALYSIS: <2–3 short sentences: what does the growing description suggest? What type of event is this? What are the key features you are noticing?>} \\
\\ 
\texttt{GUESS: <your best guess for the name of the event being described — be specific, e.g. "The French Revolution" not just "a revolution">} \\
\\ 
\texttt{Rules:} \\
\texttt{{-} You may refine your GUESS each round as more information arrives.} \\
\texttt{{-} Your GUESS should name the specific event, not just a broad category.} \\
\texttt{{-} You have up to 10 rounds total.} \\
\texttt{{-} Be concise — quality of reasoning matters more than length.} \\
\\ 
\texttt{The first sentence from the Narrator is coming now.} \\
\\ 
\\ 
\texttt{{[}Sentence 1{]}} \\
\texttt{In the year 312, a fierce confrontation erupted between the forces of Cao Shuang and Sima Yi at the capital, marking a decisive turning point in the struggle for control of the Wei state.} \\
            }
        }
    }
    & \\ \\

    \theutterance \stepcounter{utterance}  
    & & & \multicolumn{4}{p{0.6\linewidth}}{
        \cellcolor[rgb]{0.9,0.9,0.9}{
            \makecell[{{p{\linewidth}}}]{
                \texttt{\tiny{[GM$\langle$P2]}}
                \texttt{ANALYSIS: The description points to a specific military conflict in the Three Kingdoms period of China, specifically within the State of Wei. The mention of Cao Shuang and Sima Yi at the capital in 312 AD is highly specific. This appears to be a major civil war battle.} \\
\\ 
\texttt{GUESS: The Battle of Guandu (Wait, that's earlier). Let me correct: The Battle of Fancheng.} \\
            }
        }
    }
     \\ \\

    \theutterance \stepcounter{utterance}  
    & & \multicolumn{4}{p{0.6\linewidth}}{
        \cellcolor[rgb]{0.9,0.9,0.9}{
            \makecell[{{p{\linewidth}}}]{
                \texttt{\tiny{[P1$\langle$GM]}}
                \texttt{Round 1 complete. The Detective's current guess: "The Battle of Guandu (Wait, that's earlier). Let me correct: The Battle of Fancheng.". Please write a new paragraph with fresh information.} \\
            }
        }
    }
    & \\ \\

    \theutterance \stepcounter{utterance}  
    & \multicolumn{4}{p{0.6\linewidth}}{
        \cellcolor[rgb]{0.9,0.9,0.9}{
            \makecell[{{p{\linewidth}}}]{
                \texttt{\tiny{[P1$\rangle$GM]}}
                \texttt{PARAGRAPH: This pivotal clash occurred within the walls of the imperial capital, where the loyalist army under Sima Yi successfully repelled the rebel forces led by Cao Shuang, effectively ending the latter's bid to seize the throne.} \\
            }
        }
    }
    & & \\ \\

    \theutterance \stepcounter{utterance}  
    & & \multicolumn{4}{p{0.6\linewidth}}{
        \cellcolor[rgb]{0.9,0.9,0.9}{
            \makecell[{{p{\linewidth}}}]{
                \texttt{\tiny{[GM$\rangle$P2]}}
                \texttt{{[}Sentence 2{]}} \\
\texttt{This pivotal clash occurred within the walls of the imperial capital, where the loyalist army under Sima Yi successfully repelled the rebel forces led by Cao Shuang, effectively ending the latter's bid to seize the throne.} \\
            }
        }
    }
    & \\ \\

    \theutterance \stepcounter{utterance}  
    & & & \multicolumn{4}{p{0.6\linewidth}}{
        \cellcolor[rgb]{0.9,0.9,0.9}{
            \makecell[{{p{\linewidth}}}]{
                \texttt{\tiny{[GM$\langle$P2]}}
                \texttt{ANALYSIS: The details confirm a siege within the capital walls where Sima Yi defeated Cao Shuang, ending the latter's rebellion. This narrows the event to a specific defensive battle in 312 AD during the Wei state's internal conflict.} \\
\\ 
\texttt{GUESS: The Siege of Luoyang (specifically the Battle of Fancheng).} \\
            }
        }
    }
     \\ \\

    \theutterance \stepcounter{utterance}  
    & & \multicolumn{4}{p{0.6\linewidth}}{
        \cellcolor[rgb]{0.9,0.9,0.9}{
            \makecell[{{p{\linewidth}}}]{
                \texttt{\tiny{[P1$\langle$GM]}}
                \texttt{Round 2 complete. The Detective's current guess: "The Siege of Luoyang (specifically the Battle of Fancheng).". Please write a new paragraph with fresh information.} \\
            }
        }
    }
    & \\ \\

    \theutterance \stepcounter{utterance}  
    & \multicolumn{4}{p{0.6\linewidth}}{
        \cellcolor[rgb]{0.9,0.9,0.9}{
            \makecell[{{p{\linewidth}}}]{
                \texttt{\tiny{[P1$\rangle$GM]}}
                \texttt{PARAGRAPH: Following this victory, Sima Yi's troops secured the city gates and began their systematic purge of the remaining Cao Wei loyalists, a campaign that would eventually lead to the complete dissolution of the Wei regime and the consolidation of power under the Jin dynasty.} \\
            }
        }
    }
    & & \\ \\

    \theutterance \stepcounter{utterance}  
    & & \multicolumn{4}{p{0.6\linewidth}}{
        \cellcolor[rgb]{0.9,0.9,0.9}{
            \makecell[{{p{\linewidth}}}]{
                \texttt{\tiny{[GM$\rangle$P2]}}
                \texttt{{[}Sentence 3{]}} \\
\texttt{Following this victory, Sima Yi's troops secured the city gates and began their systematic purge of the remaining Cao Wei loyalists, a campaign that would eventually lead to the complete dissolution of the Wei regime and the consolidation of power under the Jin dynasty.} \\
            }
        }
    }
    & \\ \\

    \theutterance \stepcounter{utterance}  
    & & & \multicolumn{4}{p{0.6\linewidth}}{
        \cellcolor[rgb]{0.9,0.9,0.9}{
            \makecell[{{p{\linewidth}}}]{
                \texttt{\tiny{[GM$\langle$P2]}}
                \texttt{ANALYSIS: The description now highlights the political aftermath of the military victory, specifically the purge of Cao Wei loyalists and the transition of power to the Jin dynasty. This confirms the event is the culmination of Sima Yi's coup.} \\
\\ 
\texttt{GUESS: The Battle of Fancheng (312 AD).} \\
            }
        }
    }
     \\ \\

    \theutterance \stepcounter{utterance}  
    & & \multicolumn{4}{p{0.6\linewidth}}{
        \cellcolor[rgb]{0.9,0.9,0.9}{
            \makecell[{{p{\linewidth}}}]{
                \texttt{\tiny{[P1$\langle$GM]}}
                \texttt{Round 3 complete. The Detective's current guess: "The Battle of Fancheng (312 AD).". Please write a new paragraph with fresh information.} \\
            }
        }
    }
    & \\ \\

    \theutterance \stepcounter{utterance}  
    & \multicolumn{4}{p{0.6\linewidth}}{
        \cellcolor[rgb]{0.9,0.9,0.9}{
            \makecell[{{p{\linewidth}}}]{
                \texttt{\tiny{[P1$\rangle$GM]}}
                \texttt{PARAGRAPH: The aftermath of this engagement saw the immediate execution of Cao Shuang and the subsequent transfer of the imperial seal to Sima Yi, cementing his authority as the de facto ruler of the northern territories.} \\
            }
        }
    }
    & & \\ \\

    \theutterance \stepcounter{utterance}  
    & & \multicolumn{4}{p{0.6\linewidth}}{
        \cellcolor[rgb]{0.9,0.9,0.9}{
            \makecell[{{p{\linewidth}}}]{
                \texttt{\tiny{[GM$\rangle$P2]}}
                \texttt{{[}Sentence 4{]}} \\
\texttt{The aftermath of this engagement saw the immediate execution of Cao Shuang and the subsequent transfer of the imperial seal to Sima Yi, cementing his authority as the de facto ruler of the northern territories.} \\
            }
        }
    }
    & \\ \\

    \theutterance \stepcounter{utterance}  
    & & & \multicolumn{4}{p{0.6\linewidth}}{
        \cellcolor[rgb]{0.9,0.9,0.9}{
            \makecell[{{p{\linewidth}}}]{
                \texttt{\tiny{[GM$\langle$P2]}}
                \texttt{ANALYSIS: The execution of the rebel leader and the transfer of the imperial seal confirm this is the decisive moment of a coup d'état, not just a battlefield victory. The focus has shifted from the battle itself to the immediate political consolidation of power.} \\
\\ 
\texttt{GUESS: The Battle of Fancheng (312 AD).} \\
            }
        }
    }
     \\ \\

    \theutterance \stepcounter{utterance}  
    & & \multicolumn{4}{p{0.6\linewidth}}{
        \cellcolor[rgb]{0.9,0.9,0.9}{
            \makecell[{{p{\linewidth}}}]{
                \texttt{\tiny{[P1$\langle$GM]}}
                \texttt{Round 4 complete. The Detective's current guess: "The Battle of Fancheng (312 AD).". Please write a new paragraph with fresh information.} \\
            }
        }
    }
    & \\ \\

    \theutterance \stepcounter{utterance}  
    & \multicolumn{4}{p{0.6\linewidth}}{
        \cellcolor[rgb]{0.9,0.9,0.9}{
            \makecell[{{p{\linewidth}}}]{
                \texttt{\tiny{[P1$\rangle$GM]}}
                \texttt{PARAGRAPH: This decisive military engagement, fought within the fortified walls of the capital, ultimately shattered the last major resistance of the Cao Wei faction and paved the way for Sima Yi's full assumption of the regency.} \\
            }
        }
    }
    & & \\ \\

    \theutterance \stepcounter{utterance}  
    & & \multicolumn{4}{p{0.6\linewidth}}{
        \cellcolor[rgb]{0.9,0.9,0.9}{
            \makecell[{{p{\linewidth}}}]{
                \texttt{\tiny{[GM$\rangle$P2]}}
                \texttt{{[}Sentence 5{]}} \\
\texttt{This decisive military engagement, fought within the fortified walls of the capital, ultimately shattered the last major resistance of the Cao Wei faction and paved the way for Sima Yi's full assumption of the regency.} \\
            }
        }
    }
    & \\ \\

    \theutterance \stepcounter{utterance}  
    & & & \multicolumn{4}{p{0.6\linewidth}}{
        \cellcolor[rgb]{0.9,0.9,0.9}{
            \makecell[{{p{\linewidth}}}]{
                \texttt{\tiny{[GM$\langle$P2]}}
                \texttt{ANALYSIS: The description consistently emphasizes the siege within the capital walls, the defeat of Cao Shuang, and the immediate political takeover by Sima Yi. All clues point to the specific 312 AD conflict that ended the Wei faction's last stand.} \\
\\ 
\texttt{GUESS: The Battle of Fancheng (312 AD).} \\
            }
        }
    }
     \\ \\

    \theutterance \stepcounter{utterance}  
    & & \multicolumn{4}{p{0.6\linewidth}}{
        \cellcolor[rgb]{0.9,0.9,0.9}{
            \makecell[{{p{\linewidth}}}]{
                \texttt{\tiny{[P1$\langle$GM]}}
                \texttt{Round 5 complete. The Detective's current guess: "The Battle of Fancheng (312 AD).". Please write a new paragraph with fresh information.} \\
            }
        }
    }
    & \\ \\

    \theutterance \stepcounter{utterance}  
    & \multicolumn{4}{p{0.6\linewidth}}{
        \cellcolor[rgb]{0.9,0.9,0.9}{
            \makecell[{{p{\linewidth}}}]{
                \texttt{\tiny{[P1$\rangle$GM]}}
                \texttt{PARAGRAPH: The victory at this strategic stronghold not only eliminated the immediate threat to the imperial court but also triggered a chain of political maneuvers that would soon lead to the formal establishment of the Jin dynasty.} \\
            }
        }
    }
    & & \\ \\

    \theutterance \stepcounter{utterance}  
    & & \multicolumn{4}{p{0.6\linewidth}}{
        \cellcolor[rgb]{0.9,0.9,0.9}{
            \makecell[{{p{\linewidth}}}]{
                \texttt{\tiny{[GM$\rangle$P2]}}
                \texttt{{[}Sentence 6{]}} \\
\texttt{The victory at this strategic stronghold not only eliminated the immediate threat to the imperial court but also triggered a chain of political maneuvers that would soon lead to the formal establishment of the Jin dynasty.} \\
            }
        }
    }
    & \\ \\

    \theutterance \stepcounter{utterance}  
    & & & \multicolumn{4}{p{0.6\linewidth}}{
        \cellcolor[rgb]{0.9,0.9,0.9}{
            \makecell[{{p{\linewidth}}}]{
                \texttt{\tiny{[GM$\langle$P2]}}
                \texttt{ANALYSIS: The description consistently points to the 312 AD conflict where Sima Yi defeated Cao Shuang within the capital walls. The recurring themes of the siege, the purge of loyalists, and the subsequent establishment of the Jin dynasty confirm the specific historical event.} \\
\\ 
\texttt{GUESS: The Battle of Fancheng (312 AD).} \\
            }
        }
    }
     \\ \\

    \theutterance \stepcounter{utterance}  
    & & \multicolumn{4}{p{0.6\linewidth}}{
        \cellcolor[rgb]{0.9,0.9,0.9}{
            \makecell[{{p{\linewidth}}}]{
                \texttt{\tiny{[P1$\langle$GM]}}
                \texttt{Round 6 complete. The Detective's current guess: "The Battle of Fancheng (312 AD).". Please write a new paragraph with fresh information.} \\
            }
        }
    }
    & \\ \\

    \theutterance \stepcounter{utterance}  
    & \multicolumn{4}{p{0.6\linewidth}}{
        \cellcolor[rgb]{0.9,0.9,0.9}{
            \makecell[{{p{\linewidth}}}]{
                \texttt{\tiny{[P1$\rangle$GM]}}
                \texttt{PARAGRAPH: The triumph at this critical juncture allowed Sima Yi to consolidate his control over the Wei court, setting the stage for his eventual usurpation of the throne and the founding of a new imperial house.} \\
            }
        }
    }
    & & \\ \\

    \theutterance \stepcounter{utterance}  
    & & \multicolumn{4}{p{0.6\linewidth}}{
        \cellcolor[rgb]{0.9,0.9,0.9}{
            \makecell[{{p{\linewidth}}}]{
                \texttt{\tiny{[GM$\rangle$P2]}}
                \texttt{{[}Sentence 7{]}} \\
\texttt{The triumph at this critical juncture allowed Sima Yi to consolidate his control over the Wei court, setting the stage for his eventual usurpation of the throne and the founding of a new imperial house.} \\
            }
        }
    }
    & \\ \\

    \theutterance \stepcounter{utterance}  
    & & & \multicolumn{4}{p{0.6\linewidth}}{
        \cellcolor[rgb]{0.9,0.9,0.9}{
            \makecell[{{p{\linewidth}}}]{
                \texttt{\tiny{[GM$\langle$P2]}}
                \texttt{ANALYSIS: The description consistently highlights the 312 AD siege where Sima Yi defeated Cao Shuang within Luoyang's walls. The recurring themes of eliminating Wei resistance, securing the capital, and paving the way for the Jin dynasty's founding confirm the specific event.} \\
\\ 
\texttt{GUESS: The Battle of Fancheng (312 AD).} \\
            }
        }
    }
     \\ \\

    \theutterance \stepcounter{utterance}  
    & & \multicolumn{4}{p{0.6\linewidth}}{
        \cellcolor[rgb]{0.9,0.9,0.9}{
            \makecell[{{p{\linewidth}}}]{
                \texttt{\tiny{[P1$\langle$GM]}}
                \texttt{Round 7 complete. The Detective's current guess: "The Battle of Fancheng (312 AD).". Please write a new paragraph with fresh information.} \\
            }
        }
    }
    & \\ \\

    \theutterance \stepcounter{utterance}  
    & \multicolumn{4}{p{0.6\linewidth}}{
        \cellcolor[rgb]{0.9,0.9,0.9}{
            \makecell[{{p{\linewidth}}}]{
                \texttt{\tiny{[P1$\rangle$GM]}}
                \texttt{PARAGRAPH: This engagement, occurring deep within the imperial capital's defenses, served as the final major military hurdle for the Wei regime before its complete political collapse and the rise of the Jin dynasty.} \\
            }
        }
    }
    & & \\ \\

    \theutterance \stepcounter{utterance}  
    & & \multicolumn{4}{p{0.6\linewidth}}{
        \cellcolor[rgb]{0.9,0.9,0.9}{
            \makecell[{{p{\linewidth}}}]{
                \texttt{\tiny{[GM$\rangle$P2]}}
                \texttt{{[}Sentence 8{]}} \\
\texttt{This engagement, occurring deep within the imperial capital's defenses, served as the final major military hurdle for the Wei regime before its complete political collapse and the rise of the Jin dynasty.} \\
            }
        }
    }
    & \\ \\

    \theutterance \stepcounter{utterance}  
    & & & \multicolumn{4}{p{0.6\linewidth}}{
        \cellcolor[rgb]{0.9,0.9,0.9}{
            \makecell[{{p{\linewidth}}}]{
                \texttt{\tiny{[GM$\langle$P2]}}
                \texttt{ANALYSIS: The description consistently identifies a 312 AD siege within Luoyang where Sima Yi defeated Cao Shuang. The recurring details of the capital walls, the purge of Wei loyalists, and the transition to the Jin dynasty confirm the specific historical event.} \\
\\ 
\texttt{GUESS: The Battle of Fancheng (312 AD).} \\
            }
        }
    }
     \\ \\

    \theutterance \stepcounter{utterance}  
    & & \multicolumn{4}{p{0.6\linewidth}}{
        \cellcolor[rgb]{0.9,0.9,0.9}{
            \makecell[{{p{\linewidth}}}]{
                \texttt{\tiny{[P1$\langle$GM]}}
                \texttt{Round 8 complete. The Detective's current guess: "The Battle of Fancheng (312 AD).". Please write a new paragraph with fresh information.} \\
            }
        }
    }
    & \\ \\

    \theutterance \stepcounter{utterance}  
    & \multicolumn{4}{p{0.6\linewidth}}{
        \cellcolor[rgb]{0.9,0.9,0.9}{
            \makecell[{{p{\linewidth}}}]{
                \texttt{\tiny{[P1$\rangle$GM]}}
                \texttt{PARAGRAPH: The bloodshed at this fortified position within the capital directly precipitated the end of the Cao Wei dynasty, clearing the path for Sima Yi to assume the title of Regent and prepare for the establishment of the Jin.} \\
            }
        }
    }
    & & \\ \\

    \theutterance \stepcounter{utterance}  
    & & \multicolumn{4}{p{0.6\linewidth}}{
        \cellcolor[rgb]{0.9,0.9,0.9}{
            \makecell[{{p{\linewidth}}}]{
                \texttt{\tiny{[GM$\rangle$P2]}}
                \texttt{{[}Sentence 9{]}} \\
\texttt{The bloodshed at this fortified position within the capital directly precipitated the end of the Cao Wei dynasty, clearing the path for Sima Yi to assume the title of Regent and prepare for the establishment of the Jin.} \\
            }
        }
    }
    & \\ \\

    \theutterance \stepcounter{utterance}  
    & & & \multicolumn{4}{p{0.6\linewidth}}{
        \cellcolor[rgb]{0.9,0.9,0.9}{
            \makecell[{{p{\linewidth}}}]{
                \texttt{\tiny{[GM$\langle$P2]}}
                \texttt{ANALYSIS: The description consistently points to the 312 AD siege within Luoyang where Sima Yi defeated Cao Shuang. The recurring details of the capital walls, the purge of Wei loyalists, and the transition to the Jin dynasty confirm the specific historical event.} \\
\\ 
\texttt{GUESS: The Battle of Fancheng (312 AD).} \\
            }
        }
    }
     \\ \\

    \theutterance \stepcounter{utterance}  
    & & \multicolumn{4}{p{0.6\linewidth}}{
        \cellcolor[rgb]{0.9,0.9,0.9}{
            \makecell[{{p{\linewidth}}}]{
                \texttt{\tiny{[P1$\langle$GM]}}
                \texttt{Round 9 complete. The Detective's current guess: "The Battle of Fancheng (312 AD).". Please write a new paragraph with fresh information.} \\
            }
        }
    }
    & \\ \\

    \theutterance \stepcounter{utterance}  
    & \multicolumn{4}{p{0.6\linewidth}}{
        \cellcolor[rgb]{0.9,0.9,0.9}{
            \makecell[{{p{\linewidth}}}]{
                \texttt{\tiny{[P1$\rangle$GM]}}
                \texttt{PARAGRAPH: Following this confrontation, the loyalist commander Sima Yi systematically dismantled the remaining military and administrative structures of the Wei court, effectively neutralizing any organized opposition to his growing power.} \\
            }
        }
    }
    & & \\ \\

    \theutterance \stepcounter{utterance}  
    & & \multicolumn{4}{p{0.6\linewidth}}{
        \cellcolor[rgb]{0.9,0.9,0.9}{
            \makecell[{{p{\linewidth}}}]{
                \texttt{\tiny{[GM$\rangle$P2]}}
                \texttt{{[}Sentence 10{]}} \\
\texttt{Following this confrontation, the loyalist commander Sima Yi systematically dismantled the remaining military and administrative structures of the Wei court, effectively neutralizing any organized opposition to his growing power.} \\
            }
        }
    }
    & \\ \\

    \theutterance \stepcounter{utterance}  
    & & & \multicolumn{4}{p{0.6\linewidth}}{
        \cellcolor[rgb]{0.9,0.9,0.9}{
            \makecell[{{p{\linewidth}}}]{
                \texttt{\tiny{[GM$\langle$P2]}}
                \texttt{ANALYSIS: The description consistently points to the 312 AD siege within Luoyang where Sima Yi defeated Cao Shuang. The recurring details of the capital walls, the purge of Wei loyalists, and the transition to the Jin dynasty confirm the specific historical event.} \\
\\ 
\texttt{GUESS: The Battle of Fancheng (312 AD).} \\
            }
        }
    }
     \\ \\

    \theutterance \stepcounter{utterance}  
    & & & \multicolumn{2}{p{0.3\linewidth}}{
        \cellcolor[rgb]{0.9,0.9,0.9}{
            \makecell[{{p{\linewidth}}}]{
                \texttt{\tiny{[GM$|$GM]}}
                \texttt{10} \\
            }
        }
    }
    & & \\ \\

\end{supertabular}
}

\end{document}
